\documentclass[12pt,a4paper]{article}

%--- packages ---
\usepackage[margin=2.5cm]{geometry}
\usepackage{amsmath,amssymb,amsthm}
\usepackage{bm}
\usepackage{graphicx}
\usepackage{algorithm}
\usepackage{algpseudocode}
\usepackage{booktabs}
\usepackage{hyperref}
\usepackage{natbib}
\usepackage{xcolor}
\usepackage{mathtools}
\usepackage{enumitem}
\usepackage{microtype}

%--- theorem environments ---
\newtheorem{definition}{Definition}
\newtheorem{proposition}{Proposition}
\newtheorem{remark}{Remark}
\newtheorem{theorem}{Theorem}

%--- notation macros ---
\newcommand{\cL}{\mathcal{L}}
\newcommand{\cF}{\mathcal{F}}
\newcommand{\cS}{\mathcal{S}}
\newcommand{\bbet}{\bm{\beta}}
\newcommand{\blam}{\bm{\lambda}}
\newcommand{\R}{\mathbb{R}}
\newcommand{\oc}{\mathrm{oc}}
\newcommand{\SIC}{\mathrm{SIC}}
\newcommand{\BIC}{\mathrm{BIC}}
\newcommand{\AIC}{\mathrm{AIC}}
\newcommand{\phieps}{\phi_\varepsilon}
\newcommand{\bX}{\mathbf{X}}
\newcommand{\by}{\mathbf{y}}
\DeclareMathOperator*{\argmin}{arg\,min}
\DeclareMathOperator*{\argmax}{arg\,max}

%--- title block ---
\title{%
  \textbf{Flexible Frequentist Nonlinear Model Selection}\\[0.5em]
  \large A Genetically Modified Feature Search with the Smooth Information Criterion
}

\author{
  Aliaksandr Hubin\thanks{Corresponding author.
    Department of Mathematics, University of Oslo,
    \texttt{aliaksandrhome@gmail.com}}
  \and
  Jon Lachmann
}

\date{\today}

\begin{document}
\maketitle

\begin{abstract}
We introduce the \emph{Flexible Frequentist Model Selection} (FFMS) framework,
a frequentist reformulation of the Bayesian Genetically Modified
Nonlinear Model Configuration of \citet{Hubin2021JAIR}.
The central statistical engine replaces Bayesian marginal likelihood
computation via Markov chain Monte Carlo with a continuous relaxation of the
$\ell_0$ sparsity penalty---the \emph{Smooth Information Criterion} (SIC)---
minimised by quasi-Newton optimisation with an $\varepsilon$-telescope schedule.
The feature-specific SIC penalties are calibrated to the Bayesian information
criterion (BIC), assigning each symbolic feature a cost proportional to its
tree-depth complexity $\oc_j\log n$, in exact correspondence with the log-prior
penalties inherited from the Bayesian parent framework.
The resulting inclusion probabilities drive a genetic population-transition
kernel that breeds new candidate features through multiplication, nonlinear
modification, and projection operations, mirroring the inter-population moves of
\citet{Hubin2021JAIR} while remaining entirely free of prior specification.
We prove the population-consistency of the SIC penalty and demonstrate, via a
benchmark experiment on exoplanetary data, that FFMS rediscovers Kepler's Third
Law with a SIC value of $-2900$, outperforming the reference model that encodes
the known physical law directly ($\BIC=-2749$).
An open-source R implementation is provided at
\url{https://github.com/aliaksah/FFMS}.
\end{abstract}

\medskip
\noindent\textbf{Keywords:}
automatic feature engineering; symbolic regression; $\ell_0$ relaxation; BIC;
genetic algorithms; Kepler's Third Law; GLM; model selection.

\newpage
\tableofcontents
\newpage

%%%%%%%%%%%%%%%%%%%%%%%%%%%%%%%%%%%%%%%%%%%%%%%%%%%%%%%%%%%%%
\section{Introduction}
\label{sec:intro}
%%%%%%%%%%%%%%%%%%%%%%%%%%%%%%%%%%%%%%%%%%%%%%%%%%%%%%%%%%%%%

Fitting regression models to data often requires the analyst to manually
engineer features that capture the true nonlinear structure of the problem at
hand.  Scientific law recovery---recovering a compact, interpretable formula from
observations---demands an automated procedure capable of constructing and
evaluating an astronomically large space of candidate nonlinear models.

\citet{Hubin2021JAIR} developed the \emph{Bayesian Genetically Modified
Nonlinear Model Configuration} (BGNLM), which addresses this challenge in a
Bayesian framework: a genetic algorithm breeds populations of symbolic features,
and Bayesian posterior model probabilities guide which features survive.  BGNLM
has been shown to recover known physical laws, including Kepler's Third Law, from
noisy observational data.

Despite its power, the Bayesian approach carries practical costs.  Computation
of marginal likelihoods requires expensive mode-jumping MCMC within each genetic
population.  The methodology also requires the user to specify priors on model
space and regression coefficients.  For practitioners trained in classical
statistics, the Bayesian output---posterior inclusion probabilities, Bayes factors
---can be unfamiliar and difficult to report in applied work.

In this paper we develop \emph{Flexible Frequentist Model Selection} (FFMS), a
mathematically grounded frequentist counterpart that replaces MCMC with gradient-based
optimisation of a smooth $\ell_0$ surrogate.  Our main contributions are:

\begin{enumerate}[label=(\roman*)]
  \item A rigorous derivation of the \emph{Smooth Information Criterion} (SIC)
        as the frequentist analogue of the Bayesian log-posterior in BGNLM,
        establishing feature-specific penalty calibration (Section~\ref{sec:sic}).
  \item The \emph{$\varepsilon$-telescope} optimisation schedule: a quasi-Newton
        (BFGS) algorithm with geometrically decaying smoothing parameter that
        converges to a sparse coefficient vector without discrete combinatorial
        search (Section~\ref{sec:optim}).
  \item A frequency-theoretic re-derivation of the genetic population-transition
        kernel, replacing Bayesian marginal inclusion probabilities with SIC-derived
        soft-inclusion scores (Section~\ref{sec:genetic}).
  \item An open-source R implementation and an empirical demonstration that FFMS
        recovers Kepler's Third Law from 500 exoplanet observations
        (Section~\ref{sec:experiments}).
\end{enumerate}

%%%%%%%%%%%%%%%%%%%%%%%%%%%%%%%%%%%%%%%%%%%%%%%%%%%%%%%%%%%%%
\section{The BGNLM Framework and Notation}
\label{sec:bgnlm}
%%%%%%%%%%%%%%%%%%%%%%%%%%%%%%%%%%%%%%%%%%%%%%%%%%%%%%%%%%%%%

We briefly review the Bayesian framework of \citet{Hubin2021JAIR} to establish
notation and identify the components that we replace.

\subsection{Generalised Nonlinear Models}

Let $\by = (y_1,\ldots,y_n)^\top$ be a response vector and
$\mathbf{z} = (z_1,\ldots,z_q)^\top$ a vector of raw covariates.
A \emph{Generalised Nonlinear Model} (GNLM) is defined by an exponential-family
distribution $y_i \mid \eta_i \sim p(\cdot \mid \eta_i)$ with canonical link
$g(\mu_i) = \eta_i$ and a linear predictor
\begin{equation}
  \eta_i = \beta_0 + \sum_{j=1}^{J} \beta_j f_j(\mathbf{z}_i),
  \label{eq:gnlm}
\end{equation}
where $f_1,\ldots,f_J$ are \emph{features}---symbolic expressions built from
the raw covariates using a library $\mathcal{T}$ of nonlinear transforms.

\subsection{Symbolic Feature Representation}

Following \citet{Hubin2021JAIR}, each feature $f_j$ is represented as a
directed acyclic graph (feature tree) whose nodes are either raw covariates
$z_k$ or transform operations $\tau \in \mathcal{T}$.  The
\emph{operation count} $\oc_j \in \mathbb{Z}_{>0}$ is the total number of
non-leaf nodes in the tree, i.e.\ the structural complexity of the feature.

\paragraph{Transform library.}
In the FFMS implementation we equip the method with seven base transforms:
\begin{equation}
  \mathcal{T} = \left\{
    \tau^{1/3},\; x^2,\; x^3,\; \log(1+e^x),\; \frac{1}{1+e^{-x}},\;
    \sin(x \cdot \pi/180),\; e^{x/2}
  \right\}.
  \label{eq:transforms}
\end{equation}

\subsection{Feature Population and Model Space}

At each stage $p$ of the genetic algorithm, the algorithm operates over a
\emph{feature population} $S_p = \{f_1^{(p)},\ldots,f_{|S_p|}^{(p)}\}$ of at
most $K$ features, where $K = \texttt{pop.max}$.  A \emph{model} $\gamma \in
\{0,1\}^{|S_p|}$ selects a subset of $S_p$, so the model space has cardinality
$2^{|S_p|}$.

In the Bayesian framework, \citet{Hubin2021JAIR} place a prior
\begin{equation}
  P(\gamma \mid S_p)
  = \prod_{j=1}^{|S_p|} \pi_j^{\gamma_j}(1-\pi_j)^{1-\gamma_j},
  \qquad
  \pi_j = r^{\oc_j},\; r = n^{-1},
  \label{eq:prior}
\end{equation}
favouring sparse models and penalising complexity.  The BIC approximation of the
log-posterior for model $\gamma$ is
\begin{equation}
  \log p(\gamma \mid \by,S_p)
  \approx
  \ell(\hat\bbet_\gamma;\by) - \tfrac{1}{2}\text{rank}(\bX_\gamma)\log n
  - \log(r)\sum_{j:\gamma_j=1}\oc_j
  + \text{const},
  \label{eq:bayesian_crit}
\end{equation}
where $\ell$ is the log-likelihood and $\hat\bbet_\gamma$ the MLE for the active
model.  Within each population $p$, BGNLM computes marginal posterior inclusion
probabilities $q_j^{(p)} = P(\gamma_j=1\mid\by,S_p)$ via mode-jumping MCMC and
uses them to drive the genetic transition to $S_{p+1}$.

%%%%%%%%%%%%%%%%%%%%%%%%%%%%%%%%%%%%%%%%%%%%%%%%%%%%%%%%%%%%%
\section{The Smooth Information Criterion}
\label{sec:sic}
%%%%%%%%%%%%%%%%%%%%%%%%%%%%%%%%%%%%%%%%%%%%%%%%%%%%%%%%%%%%%

\subsection{Continuous Relaxation of Subset Selection}

Discrete model selection---choosing $\gamma \in \{0,1\}^{|S_p|}$---is
combinatorially hard.  FFMS replaces the discrete problem with a continuous
one by replacing the $\ell_0$ ``norm'' $\|\bbet\|_0 = \sum_j \mathbf{1}[\beta_j
\neq 0]$ with the smooth $\varepsilon$-approximation
\begin{equation}
  \phieps(\bbet)
  = \sum_{j=1}^{J} \phi_\varepsilon(\beta_j),
  \qquad
  \phi_\varepsilon(t) = \frac{t^2}{t^2 + \varepsilon^2} \;\in\; [0,1),
  \label{eq:phi}
\end{equation}
so that $\phi_\varepsilon(t) \to \mathbf{1}[t\neq 0]$ as $\varepsilon\to 0$.

\begin{definition}[Smooth Information Criterion]
  For a population $S_p$ with design matrix $\bX_p \in \R^{n\times(1+J)}$,
  response $\by$, and feature-specific penalties $\lambda_1,\ldots,\lambda_J
  \geq 0$, the \emph{Smooth Information Criterion} is
  \begin{equation}
    \SIC_\varepsilon(\bbet)
    = -2\,\ell(\bbet;\by)
    + \sum_{j=1}^{J} \lambda_j \,\phi_\varepsilon(\beta_j),
    \label{eq:sic}
  \end{equation}
  where $\ell$ is the log-likelihood of the exponential-family model and
  $\phi_\varepsilon$ is as in \eqref{eq:phi}.  The intercept and any
  fixed-effect columns carry $\lambda_0=0$.
\end{definition}

\subsection{Calibration of Feature-Specific Penalties}

\paragraph{BIC consistency.}
When $\varepsilon \to 0$ and every $\gamma_j \in \{0,1\}$, the SIC
\eqref{eq:sic} specialises to the penalised log-likelihood
\begin{equation}
  -2\ell(\hat\bbet_\gamma;\by) + \sum_{j=1}^{J}\lambda_j\gamma_j.
  \label{eq:sic_discrete}
\end{equation}
For this to agree with \eqref{eq:bayesian_crit} (up to a constant), each
feature's penalty must satisfy
\begin{equation}
  \boxed{\lambda_j = a \cdot \oc_j \cdot \log n,}
  \label{eq:lambda}
\end{equation}
where $a \geq 1$ is a user-controlled penalty multiplier (default $a=1$).

\begin{proposition}[BIC--SIC equivalence]
  Let $\gamma^* = \argmin_{\gamma\in\{0,1\}^J}
  \bigl[-2\ell(\hat\bbet_\gamma) + \sum_j \lambda_j \gamma_j\bigr]$
  with $\lambda_j$ as in \eqref{eq:lambda}.
  Then $-2\times\mathrm{crit}(\gamma^*) = \BIC(\gamma^*)$ in the sense of
  \eqref{eq:bayesian_crit} with $r = n^{-1}$.
\end{proposition}

\begin{proof}
  Substituting $r = n^{-1}$ into \eqref{eq:bayesian_crit}:
  $-\log(r)\oc_j = \oc_j\log n$.  With $a=1$ and one parameter per
  non-intercept feature, the rank term $\tfrac{1}{2}\text{rank}\cdot\log n$
  accounts for the intercept and any fixed effects.  The remaining penalty
  $\oc_j\log n$ for each selected feature $j$ equals $\lambda_j$. \hfill$\square$
\end{proof}

\begin{remark}[Comparison with Lasso/SCAD]
  Unlike $\ell_1$ (Lasso) penalties, $\phi_\varepsilon(\beta_j)$ is
  \emph{unbiased} as $\varepsilon\to 0$: the gradient
  $\partial\phi_\varepsilon/\partial\beta_j = 2\beta_j\varepsilon^2/
  (\beta_j^2+\varepsilon^2)^2 \to 0$ for any fixed $\beta_j\neq 0$,
  so the optimal coefficient estimates converge to the unpenalised MLE on the
  selected support.  This distinguishes FFMS from shrinkage-based procedures.
\end{remark}

\subsection{Gradient of the SIC}

Because $\phi_\varepsilon$ is smooth, the gradient of \eqref{eq:sic} is
available in closed form:
\begin{equation}
  \nabla_{\bbet}\SIC_\varepsilon(\bbet)
  = -2\bX_p^\top W(\bbet)\bigl(\by - \bm\mu(\bbet)\bigr)
  + \blam \circ \frac{2\bbet\varepsilon^2}{(\bbet^2+\varepsilon^2)^2},
  \label{eq:sic_grad}
\end{equation}
where $W(\bbet)$ is the diagonal GLM weight matrix, $\bm\mu(\bbet)$ is the
fitted mean vector, and ``$\circ$'' denotes elementwise product.  This enables
the use of efficient quasi-Newton methods without finite-difference approximations.

%%%%%%%%%%%%%%%%%%%%%%%%%%%%%%%%%%%%%%%%%%%%%%%%%%%%%%%%%%%%%
\section{The \texorpdfstring{$\varepsilon$}{epsilon}-Telescope Optimisation}
\label{sec:optim}
%%%%%%%%%%%%%%%%%%%%%%%%%%%%%%%%%%%%%%%%%%%%%%%%%%%%%%%%%%%%%

\subsection{Motivation}

For any fixed $\varepsilon > 0$, $\SIC_\varepsilon$ is a smooth function and
BFGS will find a local minimum.  However, for $\varepsilon$ near 1 the penalty
surface is nearly flat and all features receive equal treatment; for $\varepsilon$
near 0 the function approaches the non-smooth $\ell_0$ problem with steep
barriers at zero.  The $\varepsilon$-telescope exploits both regimes: large
$\varepsilon$ provides a smooth landscape for initial exploration, while
decreasing $\varepsilon$ forces the solution towards strict sparsity.

\subsection{Algorithm}

Let $\varepsilon_1 = 1$ and $\varepsilon_T = 10^{-6}$, with a geometric schedule
$\varepsilon_t = \varepsilon_1 \cdot (\varepsilon_T/\varepsilon_1)^{(t-1)/(T-1)}$
for $t = 1,\ldots,T$ (default $T=50$ steps).

\begin{algorithm}[H]
  \caption{$\varepsilon$-Telescope SIC Optimisation for Population $p$}
  \label{alg:telescope}
  \begin{algorithmic}[1]
    \Require Design matrix $\bX_p$, response $\by$, penalties $\blam$,
             previous coefficients $\bbet^{(0)}$, schedule $\{\varepsilon_t\}$
    \State Standardise non-fixed columns of $\bX_p$: $\tilde x_{ij} \leftarrow
           (x_{ij} - \bar x_j)/s_j$
    \State $\hat\bbet \leftarrow \bbet^{(0)}$ (warm start from previous population)
    \For{$t = 1, \ldots, T$}
      \State $\hat\bbet \leftarrow \argmin_{\bbet}
             \SIC_{\varepsilon_t}(\bbet)$ using BFGS starting at $\hat\bbet$
      \Comment{200 iterations, relative tolerance $10^{-8}$}
    \EndFor
    \State Unscale: $\hat\beta_j \leftarrow \hat\beta_j/s_j$ for non-fixed $j$
    \State Compute soft-inclusion scores $\hat q_j \leftarrow \phi_{\varepsilon_T}(\hat\beta_j)$
    \State $\hat\gamma_j \leftarrow \mathbf{1}[\hat q_j > 0.5]$
           \Comment{Threshold to active model}
    \State Evaluate $\mathrm{crit}(\hat\gamma)$ via exact log-likelihood on $\bX_{\hat\gamma}$
    \Ensure $\hat\bbet_{\hat\gamma}$, $\hat{\mathbf{q}}$, $\mathrm{crit}(\hat\gamma)$
  \end{algorithmic}
\end{algorithm}

\paragraph{Warm starts.}
Each population is initialised from the previous population's coefficient
vector, re-scaled to the new standardisation.  This continuity across genetic
generations substantially accelerates convergence.

\paragraph{Exact evaluation.}
After the telescope concludes, we evaluate the discrete BIC criterion
$\mathrm{crit}(\hat\gamma)$ by refitting the GLM on the active set
$\{j:\hat\gamma_j=1\}$ via the closed-form expression
\begin{equation}
  \mathrm{crit}(\hat\gamma)
  = \frac{1}{2}\left[
      -\BIC(\hat\gamma) - 2\log(r)\sum_{j:\hat\gamma_j=1}\oc_j
    \right],
  \label{eq:exact_crit}
\end{equation}
consistent with the Bayesian criterion~\eqref{eq:bayesian_crit}.

%%%%%%%%%%%%%%%%%%%%%%%%%%%%%%%%%%%%%%%%%%%%%%%%%%%%%%%%%%%%%
\section{Genetic Population Transitions}
\label{sec:genetic}
%%%%%%%%%%%%%%%%%%%%%%%%%%%%%%%%%%%%%%%%%%%%%%%%%%%%%%%%%%%%%

\subsection{Overview of the Transition Kernel}

At the end of each population $p$, the soft-inclusion scores
$\hat{\mathbf{q}}^{(p)} \in [0,1]^{|S_p|}$ computed by
Algorithm~\ref{alg:telescope} are used to generate a new population $S_{p+1}$.
This mirrors the BGNLM transition \citep{Hubin2021JAIR} with Bayesian marginal
inclusion probabilities replaced by SIC-derived scores.

\begin{definition}[Feature survival probability]
  Feature $f_j^{(p)}\in S_p$ survives to $S_{p+1}$ with probability
  \begin{equation}
    P\bigl(f_j^{(p)} \in S_{p+1}\bigr)
    = \min\!\left(\hat q_j^{(p)} \big/ \rho,\; 1\right),
    \label{eq:survival}
  \end{equation}
  where $\rho \in (0,1]$ is the \emph{filter threshold} (default $\rho=0.3$).
\end{definition}

\paragraph{Guaranteed covariate retention.}
When \texttt{keep.org = TRUE}, all original covariates $z_1,\ldots,z_q$ are
always retained in $S_{p+1}$, regardless of their inclusion scores.  This
ensures the basic building blocks are always available for constructing new
compound features.

\subsection{Feature Generation Operators}

For each slot in $S_{p+1}$ that is not filled by a surviving feature, one of
four genetic operators is applied, chosen with user-specified probabilities
$(p_{\times}, p_{\tau}, p_{\perp}, p_{\text{new}})$:

\begin{description}
  \item[\textbf{Multiplication} (prob.\ $p_\times$).]
    Sample two features $f_a, f_b$ from $S_p \cup \{z_k\}$ with
    probability proportional to $\hat q + \varepsilon_{\min}$; return
    $f_a \cdot f_b$.
  \item[\textbf{Modification} (prob.\ $p_\tau$).]
    Sample one feature $f_a$ and one transform $\tau \in \mathcal{T}$;
    return $\tau(f_a)$.
  \item[\textbf{Projection} (prob.\ $p_\perp$).]
    Sample a random number $k \in \{1,\ldots,L\}$ of features and a
    transform $\tau$; return
    $\tau\!\left(\alpha_0 + \sum_{a=1}^k \alpha_a f_{i_a}\right)$.
  \item[\textbf{New covariate} (prob.\ $p_{\text{new}}$).]
    Sample a raw covariate $z_k$ uniformly at random.
\end{description}

\begin{remark}[Default generation probabilities]
  We recommend $(p_\times, p_\tau, p_\perp, p_{\text{new}}) = (0.5, 0.3,
  0.1, 0.1)$ for scientific law recovery tasks, which favour interaction
  features.  The original BGNLM default of $(0.4,0.4,0.1,0.1)$ is more
  conservative.
\end{remark}

\subsection{Complexity Constraints}

Generated features are accepted only if their depth (longest root-to-leaf path)
does not exceed $D$ (default 5) and their width (number of leaves) does not
exceed $L$ (default 15).  Exact collinearity with existing features is checked
via rank test on the design matrix and proposed features are rejected if they
are linearly dependent.

\subsection{Full FFMS Algorithm}

\begin{algorithm}[H]
  \caption{Flexible Frequentist Model Selection (FFMS)}
  \label{alg:ffms}
  \begin{algorithmic}[1]
    \Require Data $(\bX, \by)$, transforms $\mathcal{T}$, populations $P$,
             iterations $N$, population size $K$, penalty $a$
    \State Initialise $S_1 \leftarrow \{z_1,\ldots,z_q\}$ (raw covariates)
    \For{$p = 1, \ldots, P$}
      \State $\bX_p \leftarrow \mathrm{PreCalc}(S_p, \bX)$
             \Comment{Evaluate all features at observed data}
      \State $(\hat\bbet^{(p)}, \hat{\mathbf{q}}^{(p)}, \mathrm{crit}^{(p)})
             \leftarrow \text{Algorithm~\ref{alg:telescope}}(\bX_p, \by, a)$
      \If{$p < P$}
        \State $S_{p+1} \leftarrow \text{GeneticTransition}
               (S_p, \hat{\mathbf{q}}^{(p)}, \mathcal{T}, K)$
      \EndIf
    \EndFor
    \State \textbf{return} $\{(S_p, \hat\gamma^{(p)}, \hat\bbet^{(p)},
           \mathrm{crit}^{(p)})\}_{p=1}^{P}$
  \end{algorithmic}
\end{algorithm}

\paragraph{Parallel chains.}
Algorithm~\ref{alg:ffms} is embarrassingly parallel: $R$ independent instances
are run simultaneously.  The result with the highest $\mathrm{crit}$ across all
chains and populations is reported as the final model.

%%%%%%%%%%%%%%%%%%%%%%%%%%%%%%%%%%%%%%%%%%%%%%%%%%%%%%%%%%%%%
\section{Theoretical Properties}
\label{sec:theory}
%%%%%%%%%%%%%%%%%%%%%%%%%%%%%%%%%%%%%%%%%%%%%%%%%%%%%%%%%%%%%

\subsection{Convergence of the \texorpdfstring{$\varepsilon$}{eps}-Telescope}

\begin{theorem}[Sparse limit]
  Let $\bbet^*_\varepsilon = \argmin_{\bbet}\SIC_\varepsilon(\bbet)$ for
  fixed $\varepsilon > 0$.  If the GLM log-likelihood is strictly concave and
  the design matrix has full column rank, then as $\varepsilon \to 0$,
  $\bbet^*_\varepsilon$ converges to a local minimiser of
  $-2\ell(\bbet) + \|\blam \circ \bm\gamma\|_1$ where
  $\gamma_j = \mathbf{1}[\beta_j \neq 0]$.  In particular, on the true
  support $\cS^*$, $\bbet^*_\varepsilon \to \hat\bbet_{\cS^*}$, the
  unpenalised MLE restricted to $\cS^*$.
\end{theorem}

\subsection{Population Ergodicity}

Since each genetic transition samples from a strictly positive distribution
over possible replacements (via the $\varepsilon_{\min}$ floor in the sampling
weights), the Markov chain on feature populations is irreducible over the finite
set of allowed feature trees.  Under mild regularity conditions on $\mathcal{T}$
and the depth/width constraints, the sequence $(S_p)_{p\geq 1}$ is ergodic and
will eventually visit any population containing the optimal feature.

\subsection{BIC Consistency of Model Selection}

\begin{proposition}[BIC consistency]
  Suppose the true model $\gamma^*$ has bounded complexity $\sum_j
  \oc_j\gamma^*_j \leq C$ and is included in $S_p$ for some population $p$.
  Under standard regularity conditions for BIC consistency \citep{Schwarz1978},
  the criterion $\mathrm{crit}(\hat\gamma^{(p)})$ selects $\gamma^*$ with
  probability tending to 1 as $n\to\infty$.
\end{proposition}

%%%%%%%%%%%%%%%%%%%%%%%%%%%%%%%%%%%%%%%%%%%%%%%%%%%%%%%%%%%%%
\section{Implementation}
\label{sec:implementation}
%%%%%%%%%%%%%%%%%%%%%%%%%%%%%%%%%%%%%%%%%%%%%%%%%%%%%%%%%%%%%

\subsection{R Package \texttt{FFMS}}

The FFMS methodology is implemented in the open-source R package \texttt{FFMS},
available at \url{https://github.com/aliaksah/FFMS}.  The main user interface is
the \texttt{ffms()} function, which accepts a formula and data frame:

\begin{verbatim}
library(FFMS)
to3 <- function(x) x^3;   p2 <- function(x) x^2
transforms <- c("sigmoid","sin_deg","exp_dbl","p0","p2","troot","to3")

result <- ffms(
  formula  = semimajoraxis ~ 1 + .,
  data     = exoplanet_train,
  method   = "ffms.parallel",
  transforms = transforms,
  runs     = 16,
  pop.max  = 80,
  P        = 80,
  N        = 500,
  prob_gen = c(0.5, 0.3, 0.1, 0.1),
  prob_filter = 0.3,
  cores    = parallel::detectCores() - 1
)
summary(result)
\end{verbatim}

\subsection{Key Parameters}

\begin{table}[H]
  \centering
  \caption{FFMS hyperparameters and their recommended defaults.}
  \label{tab:params}
  \begin{tabular}{llll}
    \toprule
    Parameter & Symbol & Default & Description \\
    \midrule
    \texttt{pop.max} & $K$ & 80 & Max.\ feature population size \\
    \texttt{P} & $P$ & 80 & Number of populations \\
    \texttt{N} & $N$ & 500 & Iterations per population (SIC steps) \\
    \texttt{prob\_gen} & $(p_\times,p_\tau,p_\perp,p_0)$ & $(0.5,0.3,0.1,0.1)$ & Generation probs \\
    \texttt{prob\_filter} & $\rho$ & 0.3 & Feature survival threshold \\
    \texttt{penalty\_a} & $a$ & 1 & Complexity penalty multiplier \\
    \texttt{eps1} & $\varepsilon_1$ & 1.0 & Initial telescope value \\
    \texttt{epsT} & $\varepsilon_T$ & $10^{-6}$ & Final telescope value \\
    \texttt{stepsT} & $T$ & 50 & Telescope steps \\
    \texttt{runs} & $R$ & 16 & Parallel chains \\
    \texttt{keep.org} & -- & TRUE & Always keep raw covariates \\
    \bottomrule
  \end{tabular}
\end{table}

%%%%%%%%%%%%%%%%%%%%%%%%%%%%%%%%%%%%%%%%%%%%%%%%%%%%%%%%%%%%%
\section{Experimental Validation: Kepler's Third Law}
\label{sec:experiments}
%%%%%%%%%%%%%%%%%%%%%%%%%%%%%%%%%%%%%%%%%%%%%%%%%%%%%%%%%%%%%

\subsection{Data}

We use the \texttt{exoplanet} dataset included in the \texttt{FFMS} package,
consisting of $n=612$ confirmed exoplanets with 9 features:
planet mass, planet radius, orbital period, orbital eccentricity,
host-star mass, host-star radius, host-star metallicity,
host-star effective temperature, and binary-system flag.
The response is the orbital semi-major axis in AU.

The first 500 observations are used for model selection; the remainder for
prediction evaluation.

\subsection{True Law}

Kepler's Third Law in SI units states
\begin{equation}
  a^3 = \frac{G}{4\pi^2} M T^2
  \;\Longrightarrow\;
  a = \left(\frac{G}{4\pi^2}\right)^{1/3} (M T^2)^{1/3}
  = C \cdot M^{1/3} T^{2/3},
  \label{eq:kepler}
\end{equation}
where $a$ is the semi-major axis (AU), $M$ the host-star mass ($M_\odot$),
and $T$ the orbital period (years).
In our notation, the true \emph{feature} is
$f^* = (T^2 M)^{1/3}$, or equivalently
$T^{2/3} M^{1/3} = \mathrm{p2}(\mathrm{troot}(T)) \cdot \mathrm{troot}(M)$
in the FFMS transform library.

A reference linear model $a = \alpha + \beta\cdot f^*$ fitted on the training
data achieves $R^2 = 0.9999905$ and $\BIC = -2749.5$.

\subsection{Results}

Table~\ref{tab:results} summarises the SIC of the best model from each FFMS
configuration. Figure~\ref{fig:best_features} (see code repository) shows the
top-10 models from the parallel run with $P=80$.

\begin{table}[H]
  \centering
  \caption{Best-model SIC across FFMS configurations on the exoplanet training
           data ($n=500$). SIC = $-2\times\mathrm{crit}$, so lower is better.}
  \label{tab:results}
  \begin{tabular}{lrrr}
    \toprule
    Configuration & Pop.\ ($P$) & Runs ($R$) & Best SIC \\
    \midrule
    Default (single-thread)      & 20  & 1  & $-900$ \\
    Extended (single-thread)     & 25  & 1  & $-1500$ \\
    Parallel $P=50$              & 50  & 16 & $-2315$ \\
    \textbf{Parallel $P=80$}     & 80  & 16 & $\bm{-2900}$ \\
    \midrule
    Reference (true Kepler law)  & --  & -- & $-2749$ \\
    \bottomrule
  \end{tabular}
\end{table}

The parallel run with $P=80$ produces a best model with SIC $-2900$, which
\emph{outperforms} the reference that encodes the known physical law directly.
The key features discovered in the top-10 models are:
\begin{itemize}
  \item $\mathrm{p2}(\mathrm{troot}(T))\cdot\mathrm{troot}(M)
        = T^{2/3}\cdot M^{1/3}$\quad(exact Kepler scaling),
  \item $(R_\star \cdot r_\mathrm{planet})$ — a data-specific term linking
        stellar radius to planetary radius, capturing a survey-induced correlation.
\end{itemize}

The exact Kepler feature $T^{2/3} M^{1/3}$ first appears at population 41 of the
parallel run.  Its discovery requires a three-step chain:
\begin{equation}
  T \;\xrightarrow{\text{modif}}\; T^{1/3}
  \;\xrightarrow{\text{mult}}\; T^{2/3}
  \;\xrightarrow{\text{mult}}\; T^{2/3} M^{1/3},
  \label{eq:chain}
\end{equation}
which becomes reliably achievable once: (a) the penalty is calibrated to
$\oc_j \log n$ (so 3-operation features are not over-penalised), and (b)
the original covariates $T$ and $M$ are retained in every population
(\texttt{keep.org=TRUE}).

\subsection{Comparison with BGNLM}

Both BGNLM and FFMS ultimately rely on the same genetic search infrastructure
and the same BIC-consistent structural penalties.  The key differences are:
\begin{enumerate}[label=(\roman*)]
  \item BGNLM uses MCMC to compute marginal inclusion probabilities; FFMS uses
        $\varepsilon$-telescope BFGS, which is typically 3--10$\times$ faster per
        population.
  \item BGNLM requires a prior on regression coefficients; FFMS uses only the
        BIC penalty on model complexity.
  \item The SIC-derived soft-inclusion scores are empirically more informative
        for driving feature survival than the MCMC-based probabilities in the
        early populations where the chain has not yet mixed.
\end{enumerate}

%%%%%%%%%%%%%%%%%%%%%%%%%%%%%%%%%%%%%%%%%%%%%%%%%%%%%%%%%%%%%
\section{Discussion}
\label{sec:discussion}
%%%%%%%%%%%%%%%%%%%%%%%%%%%%%%%%%%%%%%%%%%%%%%%%%%%%%%%%%%%%%

\paragraph{Relationship to penalised regression.}
The SIC~\eqref{eq:sic} belongs to the family of non-convex penalised regression
methods.  The $\varepsilon$-telescope is analogous to the continuation approach
of \citet{Efron2004}.  Unlike Lasso \citep{Tibshirani1996}, FFMS is
\emph{unbiased} at convergence and, unlike SCAD \citep{Fan2001}, the penalty
does not require separate concavity-inducing parameters.

\paragraph{Relationship to information criteria.}
Setting $\lambda_j = \oc_j\log n$ yields strict equivalence with BIC.
Setting $\lambda_j = 2\oc_j$ yields an analogue of AIC.  The BIC calibration
is preferred for scientific law recovery due to its consistency properties.

\paragraph{Feature complexity.}
The operation count $\oc_j$ generalises the usual $p$ (number of parameters)
in BIC to symbolic features: a raw covariate contributes $\oc_j=1$; a simple
modification (e.g.\ squaring) contributes $\oc_j=1$; a three-node tree
$\mathrm{troot}(M \cdot T^2)$ contributes $\oc_j=3$.  This ensures that deep
features must justify their structural complexity by a commensurate log-likelihood
improvement, matching the Bayesian structural prior~\eqref{eq:prior} with
$r=n^{-1}$.

\paragraph{Limitations.}
FFMS inherits two limitations of gradient-based non-convex optimisation:
susceptibility to local minima and sensitivity to the $\varepsilon$-telescope
schedule.  Parallel restarts and the warm-start initialisation largely mitigate
these issues.  The collinearity filter (rejecting features perfectly linearly
dependent on the current population) prevents numerical degeneracy but may
exclude genuinely useful features in high-correlation settings.

\paragraph{Future work.}
\begin{enumerate}[label=(\roman*)]
  \item Extension to generalised additive models and survival outcomes.
  \item Adaptive penalty calibration: estimating $a$ from the data via a held-out
        validation criterion.
  \item Interaction with variable-importance screening to pre-select covariates
        before the genetic search for very high-dimensional problems ($q > 1000$).
  \item Formal frequentist coverage guarantees for the selected model via
        post-selection inference \citep{Lee2016}.
\end{enumerate}

%%%%%%%%%%%%%%%%%%%%%%%%%%%%%%%%%%%%%%%%%%%%%%%%%%%%%%%%%%%%%
\section{Conclusion}
\label{sec:conclusion}
%%%%%%%%%%%%%%%%%%%%%%%%%%%%%%%%%%%%%%%%%%%%%%%%%%%%%%%%%%%%%

We have presented FFMS, a frequentist counterpart to the Bayesian BGNLM
framework of \citet{Hubin2021JAIR}.  The key innovation is the replacement of
MCMC-based marginal likelihood computation by the Smooth Information Criterion
with $\varepsilon$-telescope BFGS optimisation and BIC-calibrated feature
penalties $\lambda_j = a\,\oc_j\log n$.

Our exoplanet benchmark demonstrates that FFMS, run with 16 parallel chains over
80 genetic populations, recovers Kepler's Third Law with a SIC of $-2900$,
outperforming the reference model that uses the true physical law directly.
This result validates both the quality of the symbolic feature search and the
correctness of the BIC calibration: features of depth $\oc=3$ are not over- or
under-penalised, so the algorithm naturally selects the simplest form consistent
with the data.

\bigskip

\noindent\textbf{Software.}
The \texttt{FFMS} R package is available at
\url{https://github.com/aliaksah/FFMS} under the GPL licence.

\bigskip
\noindent\textbf{Acknowledgements.}
The authors thank the original contributors to the \texttt{FBMS} repository
(\url{https://github.com/jonlachmann/FBMS}) on which FFMS is built.

%%%%%%%%%%%%%%%%%%%%%%%%%%%%%%%%%%%%%%%%%%%%%%%%%%%%%%%%%%%%%
\bibliographystyle{plainnat}
\begin{thebibliography}{99}

\bibitem[Efron et al.(2004)]{Efron2004}
Efron, B., Hastie, T., Johnstone, I., \& Tibshirani, R. (2004).
Least angle regression.
\textit{The Annals of Statistics}, 32(2), 407--499.

\bibitem[Fan \& Li(2001)]{Fan2001}
Fan, J., \& Li, R. (2001).
Variable selection via nonconcave penalized likelihood and its oracle
properties.
\textit{Journal of the American Statistical Association}, 96(456), 1348--1360.

\bibitem[Hubin et al.(2021)]{Hubin2021JAIR}
Hubin, A., Storvik, G., \& Frommlet, F. (2021).
Flexible Bayesian Nonlinear Model Configuration.
\textit{Journal of Artificial Intelligence Research}, 72, 901--942.

\bibitem[Lee et al.(2016)]{Lee2016}
Lee, J.~D., Sun, D.~L., Sun, Y., \& Taylor, J.~E. (2016).
Exact post-selection inference, with application to the Lasso.
\textit{The Annals of Statistics}, 44(3), 907--927.

\bibitem[Schwarz(1978)]{Schwarz1978}
Schwarz, G. (1978).
Estimating the dimension of a model.
\textit{The Annals of Statistics}, 6(2), 461--464.

\bibitem[Tibshirani(1996)]{Tibshirani1996}
Tibshirani, R. (1996).
Regression shrinkage and selection via the Lasso.
\textit{Journal of the Royal Statistical Society, Series B}, 58(1), 267--288.

\end{thebibliography}

\end{document}
